% Блок-схема функции MainWindow::compute() — обработка нажатия кнопки "="
% ГОСТ 19.701-90, src/ui/main_window.cpp

\documentclass[a4paper]{article}
\usepackage[T2A]{fontenc}
\usepackage[utf8]{inputenc}
\usepackage[russian]{babel}
\usepackage[a4paper,margin=1.5cm]{geometry}
\usepackage{tikz}
\usetikzlibrary{shapes.geometric, arrows.meta, positioning, calc}

\tikzset{
  terminal/.style={
    rectangle, rounded corners=8pt,
    draw=black, line width=1pt,
    minimum width=4.6cm, minimum height=0.9cm,
    text centered, font=\small
  },
  process/.style={
    rectangle,
    draw=black, line width=1pt,
    minimum width=4.6cm, minimum height=0.9cm,
    text centered, font=\small
  },
  decision/.style={
    diamond, aspect=2.4,
    draw=black, line width=1pt,
    inner sep=1pt,
    minimum width=3.4cm, minimum height=0.9cm,
    text centered, font=\footnotesize
  },
  io/.style={
    trapezium, trapezium left angle=75, trapezium right angle=105,
    draw=black, line width=1pt,
    minimum width=4.6cm, minimum height=0.9cm,
    text centered, font=\small, align=center
  },
  predefined/.style={
    rectangle, double, double distance=2pt,
    draw=black, line width=1pt,
    minimum width=4.6cm, minimum height=0.9cm,
    text centered, font=\small
  },
  arrow/.style={->, >=Stealth, line width=1pt}
}

\begin{document}
\begin{center}
\textbf{Блок-схема функции \texttt{MainWindow::compute()}}
\end{center}

\begin{center}
\begin{tikzpicture}[node distance=0.7cm]

\node (start)  [terminal]                          {Начало};
\node (read)   [io,        below=of start]         {Прочитать строку из поля ввода};
\node (trim)   [process,   below=of read]          {expr = строка без пробелов по краям};
\node (empty)  [decision,  below=of trim]          {expr пустая?};

\node (clrRes) [io,        right=2.5cm of empty]   {Очистить поле результата};
\node (retE)   [terminal,  below=of clrRes]        {Конец};

\node (calc)   [predefined, below=of empty]        {r = evaluate(expr)};

\node (err)    [decision,  below=of calc]          {Произошла ошибка?};

\node (histE)  [process,   below=of err, xshift=-3.4cm]
                                                   {history = expr};
\node (showE)  [io,        below=of histE, align=center]
                                                   {Показать сообщение об ошибке (красным)};

\node (histOk) [process,   below=of err, xshift=3.4cm]
                                                   {history = expr + " ="};
\node (fmt)    [predefined, below=of histOk]       {text = format\_result(r)};
\node (showOk) [io,        below=of fmt, align=center]
                                                   {Показать text (зелёным)};

\node (finish) [terminal,  below=2.0cm of err, yshift=-3.6cm]
                                                   {Конец};

\draw [arrow] (start) -- (read);
\draw [arrow] (read)  -- (trim);
\draw [arrow] (trim)  -- (empty);

\draw [arrow] (empty) -- node[above] {Да} (clrRes);
\draw [arrow] (clrRes) -- (retE);

\draw [arrow] (empty) -- node[left]  {Нет} (calc);
\draw [arrow] (calc)  -- (err);

\draw [arrow] (err.west)  -| node[above, pos=0.25] {Да}  (histE.north);
\draw [arrow] (err.east)  -| node[above, pos=0.25] {Нет} (histOk.north);

\draw [arrow] (histE)  -- (showE);
\draw [arrow] (histOk) -- (fmt);
\draw [arrow] (fmt)    -- (showOk);

\draw [arrow] (showE.south)  |- (finish.west);
\draw [arrow] (showOk.south) |- (finish.east);

\end{tikzpicture}
\end{center}

\end{document}
